\documentclass[12pt]{article}
\usepackage[T1]{fontenc}
\usepackage[utf8]{inputenc}
\usepackage{lmodern}
\usepackage{textcomp}
\usepackage{enumitem}
\usepackage{geometry}
\geometry{margin=1in}
\begin{document}
\begin{center}
Answer Key - History - Graduate Level Assessment 7 \
\small (Answers are ordered exactly as in the question paper.)
\end{center}

\section*{Multiple Choice}
\begin{enumerate}[label=\arabic*.]
\item \textbf{Answer:} C
\\ \textit{Options:} A) 3100 B.P.; B) 4100 B.P.; C) 5100 B.P.; D) 6100 B.P.
\\ \textit{Note:} The correct answer of 5100 B.P. reflects the emergence of the first pharaohs in Egypt around 3100 BCE, marking the unification of Upper and Lower Egypt and the beginning of the Early Dynastic Period, which is a pivotal moment in ancient Egyptian history.
\item \textbf{Answer:} A
\\ \textit{Options:} A) establishment of permanent settlements; B) development of agriculture; C) domestication of animals; D) creation of written language; E) symbolic expression through the production of art
\\ \textit{Note:} The establishment of permanent settlements is the most ancient characteristic of modern humans because it reflects the transition from nomadic lifestyles to sedentism, which began with the advent of agriculture and the Neolithic Revolution, predating the cultural developments of the Late Stone Age and Upper Paleolithic.
\item \textbf{Answer:} B
\\ \textit{Options:} A) the emergence of complex social hierarchies and political systems; B) a growth in cultural diversity; C) the formation of a single global civilization; D) the onset of a new ice age; E) the introduction of farming from neighboring continents
\\ \textit{Note:} During the Holocene epoch, Australia, like other regions of the world, experienced significant advancements in human culture, leading to increased cultural diversity as societies adapted to changing environments and developed new technologies and social structures.
\item \textbf{Answer:} D
\\ \textit{Options:} A) both a and b; B) Clovis; C) Gravettian; D) Mousterian; E) Acheulean
\\ \textit{Note:} The Mousterian tool technology, characterized by the production of flake tools and associated with the Middle Paleolithic period, is specifically linked to Neandertals, reflecting their advanced craftsmanship and adaptation to their environment.
\item \textbf{Answer:} C
\\ \textit{Options:} A) a diet high in calcium or low in calcium.; B) a diet high in saturated fats or low in saturated fats.; C) mostly grains or mostly nuts and fruits.; D) mostly root vegetables or mostly leafy greens.; E) a diet rich in sugar or low in sugar.
\\ \textit{Note:} The correct answer is based on the fact that different plants, such as grains and nuts/fruits, have distinct carbon isotopic signatures due to their varying photosynthetic pathways, allowing researchers to determine an individual's dietary preferences through the analysis of 13 C levels in their bones.
\end{enumerate}

\section*{True / False}
\begin{enumerate}[label=\arabic*.]
\setcounter{enumi}{5}
\item \textbf{Answer:} False
\\ \textit{Options:} A) True; B) False
\\ \textit{Note:} The correct answer is: seizing prey, biting off pieces of vegetation, or grasping dead and decaying matter
\item \textbf{Answer:} False
\\ \textit{Options:} A) True; B) False
\\ \textit{Note:} The correct answer is: 0.3
\item \textbf{Answer:} False
\\ \textit{Options:} A) True; B) False
\\ \textit{Note:} The correct answer is: 1914
\item \textbf{Answer:} False
\\ \textit{Options:} A) True; B) False
\\ \textit{Note:} The correct answer is: dance
\item \textbf{Answer:} True
\\ \textit{Options:} A) True; B) False
\\ \textit{Note:} According to the passage: commercialism
\end{enumerate}

\section*{Long Form Response}
\begin{enumerate}[label=\arabic*.]
\setcounter{enumi}{10}
\item \textbf{Answer:} highest
\\ \textit{Note:} Expected answer: highest
\item \textbf{Answer:} Rajasthan
\\ \textit{Note:} Expected answer: Rajasthan
\end{enumerate}

\vfill
\noindent\textit{End of Answer Key}
\end{document}